% ---------------------------------------------------------------------------
% fuzzyclust book — shared preamble and macros
% Every chapter uses ONLY these macros for algorithms, code and figures so
% that the book stays uniform. See book/README.md for the conventions.
% Visual identity (colors, display face, logo): cover/brand.tex.
% ---------------------------------------------------------------------------
\usepackage[utf8]{inputenc}
\usepackage[T1]{fontenc}
\usepackage[english]{babel}
% --- type: Libertinus Serif and Libertinus Math for text and formulas,
%     Raleway for display, IBM Plex Mono for code ----------------------------
\usepackage{amsmath,amsthm}
\usepackage[lining]{libertinus}
\usepackage{libertinust1math}
\usepackage{bm,mathtools}
\usepackage[type1,lining]{raleway}
\usepackage[scale=0.86]{plex-mono}
\usepackage[letterspace=120]{microtype}
\usepackage{graphicx}
\usepackage{xcolor}
\usepackage{booktabs,tabularx,array,multirow,longtable,float}
\usepackage{enumitem}
\usepackage{caption,subcaption}
\usepackage{tikz}
\usetikzlibrary{arrows.meta,positioning,shapes.geometric,calc,fit,backgrounds,decorations.pathreplacing,matrix,fadings}
\usepackage{pgfplots}
\pgfplotsset{compat=1.18}
\usepackage[plain,vlined,linesnumbered]{algorithm2e}
\usepackage{listings}
\usepackage[most]{tcolorbox}
\usepackage{pdfpages}
\usepackage{titlesec,titletoc}
\usepackage{fancyhdr,emptypage}
\usepackage{needspace}
\usepackage{makeidx}
% ---------------------------------------------------------------------------
% fuzzyclust visual identity: colors, display font and the logo mark.
% Shared by the book (preamble.tex), the cover, the banner and the software
% documentation. Colors: Pannon blue for structure, the Okabe--Ito triple
% (blue, vermilion, bluish green) of the logo and of every data figure.
% ---------------------------------------------------------------------------
\definecolor{fcnavy}{HTML}{14304A}   % cover and dark backgrounds
\definecolor{fcbrand}{HTML}{285780}  % headings, rules, links
\definecolor{fcink}{HTML}{1F2A36}    % running text in headings
\definecolor{fcblue}{HTML}{0072B2}
\definecolor{fcverm}{HTML}{D55E00}
\definecolor{fcgreen}{HTML}{009E73}
\definecolor{fcpurple}{HTML}{CC79A7}
\definecolor{fcorange}{HTML}{E69F00}
\definecolor{fcsky}{HTML}{56B4E9}
\definecolor{fcgray}{HTML}{6B7785}
\definecolor{fcpaper}{HTML}{F4F6F8}
\definecolor{fclight}{HTML}{9CC3E6}  % light accent on dark backgrounds

% display face: Raleway (Type 1 via the raleway package)
\providecommand{\fcdisplay}{\fontfamily{Raleway-TLF}\selectfont}
\providecommand{\fcblack}{\fontfamily{Raleway-TLF}\fontseries{k}\selectfont}
\providecommand{\fcbold}{\fontfamily{Raleway-TLF}\fontseries{b}\selectfont}
\providecommand{\fcsemi}{\fontfamily{Raleway-TLF}\fontseries{sb}\selectfont}
\providecommand{\fcmedium}{\fontfamily{Raleway-TLF}\fontseries{mb}\selectfont}
\providecommand{\fclightface}{\fontfamily{Raleway-TLF}\fontseries{l}\selectfont}

% logo mark: three membership clouds; \fclogo[<scale>]{<opacity factor>}
% drawn with its center at the origin, unit = 1cm at scale 1
\newcommand{\fclogocloud}[3]{% x, y, color
  \foreach \r/\o in {1.25/0.18,1.0/0.32,0.68/0.55,0.36/0.9}{%
    \fill[#3, fill opacity=\o] (#1,#2) circle (\r);}}
\newcommand{\fclogo}[1][1]{%
  \begin{tikzpicture}[scale=#1, every node/.style={transform shape}]
    \fclogocloud{-0.9}{-0.52}{fcblue}
    \fclogocloud{0.9}{-0.52}{fcverm}
    \fclogocloud{0}{1.04}{fcgreen}
  \end{tikzpicture}}

\usepackage[colorlinks=true,linktoc=page,linkcolor=fcbrand,citecolor=fcgreen!80!black,urlcolor=fcbrand]{hyperref}
\usepackage{cleveref}
\usepackage{csquotes}
\usepackage[backend=biber,style=numeric,sorting=none,sortcites=true,maxbibnames=6,doi=true,url=false,isbn=false,eprint=false]{biblatex}
\usepackage{imakeidx}
\makeindex[columns=2]
\renewcommand*{\bibfont}{\small\raggedright}
\setcounter{biburlnumpenalty}{100}
\setcounter{biburllcpenalty}{100}
\setcounter{biburlucpenalty}{100}

% --- page layout ---------------------------------------------------------------
% text block of the standard book class (360 pt measure, which the TikZ figures
% are drawn for); raggedbottom keeps the interline spacing constant
\raggedbottom
\widowpenalty=10000
\clubpenalty=10000
\displaywidowpenalty=3000
\setlength{\parskip}{3pt plus 1pt}
\setlength{\emergencystretch}{1.5em}

% --- chapter and section headings ---------------------------------------------
\titleformat{\chapter}[display]
  {\raggedright\color{fcink}}
  {\fcsemi\footnotesize\color{fcgray}\textls[200]{\MakeUppercase{\chaptertitlename}}\\[2pt]%
   {\fcblack\fontsize{58}{58}\selectfont\color{fcbrand}\thechapter}\hfill\raisebox{2pt}{\fclogo[0.21]}}
  {10pt}
  {\fcbold\fontsize{25}{30}\selectfont}
  [\vspace{10pt}{\color{fcverm}\rule[0.25pt]{1.6cm}{2.2pt}}{\color{fcbrand!45}\rule[1.2pt]{\dimexpr\linewidth-1.6cm\relax}{0.5pt}}]
\titleformat{name=\chapter,numberless}[display]
  {\raggedright\color{fcink}}
  {}
  {0pt}
  {\fcbold\fontsize{25}{30}\selectfont}
  [\vspace{10pt}{\color{fcverm}\rule[0.25pt]{1.6cm}{2.2pt}}{\color{fcbrand!45}\rule[1.2pt]{\dimexpr\linewidth-1.6cm\relax}{0.5pt}}]
\titlespacing*{\chapter}{0pt}{0pt}{30pt}
\titleformat{\section}
  {\fcbold\Large\color{fcink}\raggedright}
  {\color{fcbrand}\thesection}{0.75em}{}
\titleformat{\subsection}
  {\fcbold\large\color{fcink}\raggedright}
  {\color{fcbrand}\thesubsection}{0.7em}{}
\titleformat{\subsubsection}
  {\fcsemi\normalsize\color{fcink}\raggedright}
  {\thesubsubsection}{0.6em}{}
\titleformat{\paragraph}[runin]
  {\fcsemi\normalsize\color{fcink}}{}{0pt}{}
\titlespacing*{\section}{0pt}{18pt plus 4pt minus 2pt}{8pt plus 2pt}
\titlespacing*{\subsection}{0pt}{14pt plus 3pt minus 2pt}{6pt plus 2pt}
\titlespacing*{\subsubsection}{0pt}{10pt plus 3pt minus 2pt}{4pt plus 1pt}
\titlespacing*{\paragraph}{0pt}{6pt plus 2pt minus 1pt}{0.6em}
% a heading keeps at least three lines of text below it on its page (\Needspace
% measures the page exactly; \needspace misbehaves with \raggedbottom)
\preto\section{\Needspace{5\baselineskip}}
\preto\subsection{\Needspace{4\baselineskip}}

% --- table of contents and lists of figures and tables -------------------------
\setcounter{tocdepth}{1}
\contentsmargin{2.55em}
\titlecontents{chapter}[2.2em]
  {\addvspace{9pt}\fcbold\color{fcink}}
  {\contentslabel[\color{fcbrand}\thecontentslabel]{2.2em}}
  {\hspace*{-2.2em}}
  {\hfill\color{fcbrand}\contentspage}
\titlecontents{section}[5.0em]
  {\small}
  {\contentslabel[\color{fcgray}\thecontentslabel]{2.8em}}
  {\hspace*{-2.8em}}
  {\titlerule*[0.7pc]{\color{fcgray!60}.}\contentspage}
\titlecontents{figure}[2.8em]
  {\small}
  {\contentslabel[\color{fcbrand}\thecontentslabel]{2.8em}}
  {\hspace*{-2.8em}}
  {\titlerule*[0.7pc]{\color{fcgray!60}.}\contentspage}
\titlecontents{table}[2.8em]
  {\small}
  {\contentslabel[\color{fcbrand}\thecontentslabel]{2.8em}}
  {\hspace*{-2.8em}}
  {\titlerule*[0.7pc]{\color{fcgray!60}.}\contentspage}

% --- running heads -------------------------------------------------------------
\pagestyle{fancy}
\fancyhf{}
\fancyhead[LE]{\fcsemi\footnotesize{\color{fcbrand}\thepage}\hspace{1.2em}{\color{fcgray}\leftmark}}
\fancyhead[RO]{\fcsemi\footnotesize{\color{fcgray}\rightmark}\hspace{1.2em}{\color{fcbrand}\thepage}}
\renewcommand{\headrulewidth}{0pt}
\makeatletter
\renewcommand{\chaptermark}[1]{\markboth{\ifnum\c@secnumdepth>\m@ne\if@mainmatter\@chapapp\ \thechapter\hspace{0.6em}\fi\fi #1}{\ifnum\c@secnumdepth>\m@ne\if@mainmatter\@chapapp\ \thechapter\hspace{0.6em}\fi\fi #1}}
\makeatother
\renewcommand{\sectionmark}[1]{\markright{\thesection\hspace{0.6em}#1}}
\fancypagestyle{plain}{\fancyhf{}\fancyfoot[C]{\fcsemi\footnotesize\color{fcgray}\thepage}}

% --- captions ------------------------------------------------------------------
\DeclareCaptionFont{fclabel}{\fcbold\color{fcbrand}}
\captionsetup{font=small,labelfont=fclabel,labelsep=period,skip=6pt}
\captionsetup[table]{position=top}

% --- notation macros ---------------------------------------------------------
\newcommand{\X}{\mathbf{X}}
\newcommand{\xk}{\mathbf{x}_k}
\newcommand{\vi}{\mathbf{v}_i}
\newcommand{\V}{\mathbf{V}}
\newcommand{\U}{\mathbf{U}}
\newcommand{\mik}{\mu_{ik}}
\newcommand{\A}{\mathbf{A}}
\newcommand{\Ai}{\mathbf{A}_i}
\newcommand{\Fi}{\mathbf{F}_i}
\newcommand{\Dik}{D_{ik}}
\newcommand{\Dsq}{D^2_{ik}}
\newcommand{\J}{J}
\newcommand{\R}{\mathbb{R}}
\newcommand{\T}{^{\mathsf{T}}}
\newcommand{\norm}[1]{\left\lVert #1 \right\rVert}
\newcommand{\abs}[1]{\left\lvert #1 \right\rvert}
\DeclareMathOperator{\argmin}{arg\,min}
\DeclareMathOperator{\argmax}{arg\,max}
\DeclareMathOperator{\diag}{diag}
\DeclareMathOperator{\tr}{tr}
% \code: typewriter identifiers; a line may break after an underscore (long
% snake_case names would otherwise produce overfull lines in the narrow measure)
\makeatletter
% break after an underscore only when more of the name follows (a trailing
% underscore, as in cluster_centers_, must not be followed by a line break)
\def\fc@us{\textunderscore\futurelet\fc@next\fc@usbreak}
\def\fc@usbreak{\ifx\fc@next\egroup\else\allowbreak\fi}
\DeclareRobustCommand{\code}[1]{{\ttfamily\let\_\fc@us#1}}
\makeatother
\pdfstringdefDisableCommands{\def\code#1{#1}} % plain text in PDF bookmarks
\newcommand{\py}[1]{\lstinline[language=Python]|#1|}

% --- theorem-like environments ----------------------------------------------
\newtheoremstyle{fcdef}{8pt}{8pt}{\normalfont}{}{\fcbold\color{fcbrand}}{.}{0.5em}{}
\newtheoremstyle{fcrem}{8pt}{8pt}{\normalfont}{}{\fcsemi\color{fcink}}{.}{0.5em}{}
\theoremstyle{fcdef}
\newtheorem{definition}{Definition}[chapter]
\newtheorem{example}{Example}[chapter]
\theoremstyle{fcrem}
\newtheorem{remark}{Remark}[chapter]

% --- listings ----------------------------------------------------------------
\lstdefinestyle{fcpython}{
  language=Python,
  basicstyle=\ttfamily\footnotesize,
  keywordstyle=\color{fcbrand}\bfseries,
  stringstyle=\color{fcgreen!85!black},
  commentstyle=\color{fcgray}\itshape,
  showstringspaces=false,
  breaklines=true,
  breakatwhitespace=true,
  columns=fullflexible,
  keepspaces=true,
  frame=leftline,
  framerule=1.6pt,
  rulecolor=\color{fcbrand!65},
  framesep=6pt,
  backgroundcolor=\color{fcpaper},
  xleftmargin=8pt, xrightmargin=0pt,
  aboveskip=8pt, belowskip=8pt,
  rangebeginprefix=\#\ [snippet:,
  rangebeginsuffix=],
  rangeendprefix=\#\ [/snippet:,
  rangeendsuffix=],
  includerangemarker=false,
  morekeywords={self,None,True,False},
}
\lstset{style=fcpython}

% \pysnippet{examples/01_clustering/plot_01_fcm_basics.py}{fit}
%   -> code between "# [snippet:fit]" and "# [/snippet:fit]" in that script
% a listing never starts with fewer than four lines at the foot of a page
\newcommand{\pysnippet}[2]{\Needspace{5\baselineskip}\lstinputlisting[linerange=#2-#2]{../#1}}
% \pyfile{examples/...py} -> whole script (used in the example appendix)
\newcommand{\pyfile}[1]{\Needspace{5\baselineskip}\lstinputlisting{../#1}}
\AtBeginEnvironment{lstlisting}{\Needspace{5\baselineskip}}

% --- figures from the example scripts ----------------------------------------
% \exfig[width]{01_clustering/plot_03_gk_ellipses_1}{caption}{label}
% Files are produced by book/build_figures.py into book/figures/<section>/<script>_<n>.pdf
% \exfig[width]{path}[short caption]{caption}{label}; the short caption goes to the list of figures
\NewDocumentCommand{\exfig}{O{\textwidth} m o m m}{%
  \begin{figure}[htbp]\centering
    \includegraphics[width=#1]{figures/#2.pdf}
    \IfNoValueTF{#3}{\caption{#4}}{\caption[{#3}]{#4}}\label{fig:#5}
  \end{figure}}
\renewcommand{\topfraction}{0.85}
\renewcommand{\bottomfraction}{0.6}
\renewcommand{\textfraction}{0.12}
\renewcommand{\floatpagefraction}{0.75}
\setcounter{topnumber}{2}
\setcounter{totalnumber}{3}

% --- boxes --------------------------------------------------------------------
% One geometry for every box: a colored rule on the left, a light tint, the
% title in the display face. Only the color changes with the kind of box.
\tcbset{
  fcbase/.style={enhanced, breakable, sharp corners, frame hidden, boxrule=0pt,
    borderline west={2.2pt}{0pt}{#1}, colback=#1!6!white,
    left=9pt, right=8pt, top=6pt, bottom=6pt, before skip=10pt, after skip=10pt,
    detach title, before upper={\tcbtitle\par\vspace{3pt}}, lines before break=4,
    fonttitle=\fcbold\small, coltitle=#1!75!black},
}
% algorithm box: \begin{algobox}{Title}{label} <algorithm2e body> \end{algobox}
% (a float: an algorithm that does not fit moves to the next page and the text
% continues instead of leaving the foot of the page empty)
% numbered per chapter, so that \cref{alg:...} prints "algorithm 3.2"
\newtcolorbox[auto counter, number within=chapter, crefname={algorithm}{algorithms},
  Crefname={Algorithm}{Algorithms}]{algobox}[2]{fcbase=fcbrand, unbreakable, float=tbp,
  title={Algorithm \thetcbcounter\enspace\textcolor{fcink}{#1}}, label={alg:#2}}
% summary box at the end of a theory section
% (a float like the algorithm boxes, in whose order it stays: the text flows on
% instead of leaving the foot of a page empty)
\newtcolorbox{summarybox}[1]{fcbase=fcgreen, unbreakable, float=tbp,
  title={In a nutshell\enspace\textcolor{fcink}{#1}}}
% practical hints
\newtcolorbox{hintbox}[1]{fcbase=fcorange, unbreakable,
  title={#1}}
% roadmap at the beginning of every chapter
\newtcolorbox{roadmap}{enhanced, sharp corners, frame hidden, boxrule=0pt,
  borderline west={2.2pt}{0pt}{fcverm}, colback=fcpaper,
  left=10pt, right=10pt, top=7pt, bottom=7pt, before skip=0pt, after skip=14pt,
  fontupper=\small\color{fcink}}
% API box: parameters of an estimator
\newtcolorbox{apibox}[1]{fcbase=fcsky!80!fcbrand,
  title={\code{#1}}, fonttitle=\small\bfseries\ttfamily,
  before upper={\tcbtitle\par\vspace{3pt}\emergencystretch=2.5em\hbadness=10000}}

% --- TikZ styles for concept figures -----------------------------------------
\tikzset{
  block/.style={draw=fcbrand, fill=fcbrand!7, rounded corners=2pt, thick, minimum height=8mm, minimum width=22mm, align=center, font=\small},
  data/.style={draw=fcgreen!70!black, fill=fcgreen!8, rounded corners=2pt, thick, minimum height=8mm, align=center, font=\small},
  result/.style={draw=fcverm!70!black, fill=fcverm!8, rounded corners=2pt, thick, minimum height=8mm, align=center, font=\small},
  note/.style={font=\footnotesize\itshape, text=fcgray, align=center},
  arrow/.style={-{Latex[length=2mm]}, thick, fcgray!80!black},
  every node/.append style={execute at begin node={\hyphenpenalty=10000\exhyphenpenalty=10000}},
}

\graphicspath{{figures/}}
\captionsetup{font=small}
\setlist{nosep}
